\section{Toward Polynomially Hard Obfuscation} \label{sec-poly}

The subexponential assumption in Theorem~\ref{thm-adaptive} comes from two guessing steps. Each signing hop guesses the adaptive message so that the reduction can puncture at the two points used in that hop. The trigger-forgery reduction guesses the complete forgery slot. This section removes the signing-hop guesses, so that the loss from adaptivity is confined to the trigger-forgery analysis. We prove the signing hops exactly, with no loss, and we isolate the remaining step as a precisely stated requirement on a message-binding mechanism. We do not resolve that requirement, and we do not claim a theorem under polynomially hard iO. We present the construction and the exact signing lemma because the lemma is clean and, in our view, the right first step, and because the remaining requirement is a sharp and self-contained question.

The scheme is the salted scheme of Section~\ref{sec-salted}, with a longer salt. The salt is no longer inert random bits. It is a one-time encoding of the trigger data, so that a single signature contains its own programmed point inside its salt. We first fix the encoding.

\subsection{The Encoded-Salt Construction} \label{sec-poly-constr}

Let $\CRHF \colon \bits^{\msglen} \to \bits^{\hlen}$ be a collision-resistant hash function with $\hlen = 2\secp$. Let $\PRF_5$ be a puncturable PRF with domain $\bits^{\secp}$, the space of \emph{handles}, and range $\bits^{\plen}$, where $\plen$ is the length of the payload defined next. The \emph{payload} is a tuple $(c, z, h, t) \in \Zp \times \Zp \times \bits^{\hlen} \times \bits^{\secp}$, and $\plen$ is its length in bits; its last component $t$ is an authentication tag. Let $\PRF_6$ be a puncturable PRF with range $\bits^{\secp}$, which computes the tag. A salt is a pair $(\handle, w) \in \bits^{\secp} \times \bits^{\plen}$.

\begin{figure}[t]
\begin{center}
\fbox{
\begin{minipage}{4.7in}
\begin{tabbing}
12\=12\=12\=\kill
\textbf{Circuit} $\CircP[X, k_1, k_2, k_5, k_6](R, m, (\handle, w))$: \\
\> $\pad \gets \PRFEv_5(k_5, \handle)$ \comment{one-time pad for this handle} \\
\> $(c, z, h, t) \gets w \oplus \pad$ \comment{decode the payload} \\
\> If $t = \PRFEv_6(k_6, (\handle, c, z, h))$ \textbf{and} $h = \CRHF(m)$ \textbf{and} $R = g^{z} X^{-c}$ \textbf{and} $c \ne 0$: \\
\> \> Return $c$ \comment{trigger branch} \\
\> $a \gets \PRFEv_1(k_1, (R, m, \handle, w)) \;;\; b \gets \PRFEv_2(k_2, (R, m, \handle, w))$ \\
\> Return $\big(f(R \cdot X^{a} \cdot g^{b}) + a\big) \bmod p$ \comment{main branch}
\end{tabbing}
\end{minipage}
}
\end{center}
\vspace{-6pt}
\caption{The encoded-salt hash circuit. It is padded to a fixed polynomial size before obfuscation.}\label{fig-circuit-poly}
\vspace{-3pt}\hrulefill
\vspace{-1ex}
\end{figure}

\begin{construction} \label{constr-poly}
The salted keyed hash family $\HF^{\mathsf{poly}} = (\HFKg, \HFEv)$ is defined as follows. Algorithm $\HFKg$ on input $X$ samples keys $k_1, k_2, k_5, k_6$ for the four PRFs and outputs $\hk \gets \iO(\usecp, \CircP[X, k_1, k_2, k_5, k_6])$. Algorithm $\HFEv$ evaluates the obfuscated circuit. The scheme is $\SchSScheme[\GG, \HF^{\mathsf{poly}}]$, the salted Schnorr scheme of Construction~\ref{constr-salted-sch}.
\end{construction}

\paragraph*{Discussion.} The circuit has the same two branches as Figure~\ref{fig-circuit-salted}. The difference is how the trigger branch decides to activate. It no longer derives the trigger point from the slot by a PRF. Instead it decodes a payload from the salt, using a one-time pad keyed by the handle, and activates if the decoded payload is a well-formed, authenticated, message-bound programming instruction. The four checks are, in order: the tag $t$ authenticates the triple $(c, z, h)$ under the handle, the field $h$ equals the hash of the message, the point $R$ is the programmed point $g^{z} X^{-c}$, and $c$ is nonzero (as in Section~\ref{sec-io-uf}, so that a trigger forgery is excluded from the CDL solution set). An honest signer draws a uniform salt. Then the decoded payload is a uniform string, the tag check fails except with probability $2^{-\secp}$, and the trigger does not activate, so honest signing uses the main branch and is unchanged. The reduction signs an adaptively chosen $m$ by choosing $(c, z)$ uniformly, computing $h = \CRHF(m)$ and the tag $t = \PRFEv_6(k_6, (\handle, c, z, h))$ for a fresh handle $\handle$, and encrypting the payload under the pad $\PRFEv_5(k_5, \handle)$. It outputs the signature $(g^{z} X^{-c}, z, (\handle, w))$. No secret key is used, and the pad is evaluated at the handle, which is the reduction's own fresh randomness and does not contain the message.

\subsection{The Signing Hops Are Exact} \label{sec-poly-lemma}

We isolate the gain of the encoding. Consider the two ways of answering one signing query on a message $m$. The \emph{honest procedure} draws $r \getsr \Zp$ and a uniform salt $\salt$, and returns $\big(g^r,\ (r + \HFEv(\hk, (g^r, m, \salt)) \cdot x) \bmod p,\ \salt\big)$. The \emph{trigger procedure} draws $c, z \getsr \Zp$ and a fresh handle $\handle$, sets $h \gets \CRHF(m)$, $t \gets \PRFEv_6(k_6, (\handle, c, z, h))$, $\pad \gets \PRFEv_5(k_5, \handle)$, $w \gets (c, z, h, t) \oplus \pad$, and returns $(g^{z} X^{-c}, z, (\handle, w))$.

\begin{lemma} \label{lem-exact-sign}
Fix a query index $j$. Consider two games that are identical except that query $j$ is answered by the honest procedure in one and by the trigger procedure in the other, and in which, in both games, the handle used for query $j$ is a value $\handle_j$ sampled uniformly at key-generation time, before any query, and used only for query $j$. Then the two games are indistinguishable, with the difference bounded by the iO advantage plus the puncturable-PRF advantage against $\PRF_5$ at the single point $\handle_j$ plus $2^{-\secp}$.
\end{lemma}

\begin{proof}
The one point that needs care is that the reduction must know the handle $\handle_j$ when it builds the obfuscated circuit, which happens at key generation, before query $j$ is made. This is possible precisely because the handle is message-independent: the reduction samples $\handle_j$ at key generation and commits to using it for query $j$. Since $\handle_j$ is uniform whether it is sampled in advance or at query time, this changes no distribution; it only lets the reduction hardwire.

We first replace the pad of query $j$ by a uniform value. The pad is $\pad = \PRFEv_5(k_5, \handle_j)$. At key generation, puncture $k_5$ at the single point $\handle_j$ and hardwire the value $\pad$ into the circuit at that point; since the hardwired value equals the real value, the punctured circuit computes the same function as the original, so the obfuscations are indistinguishable, at the cost of the iO advantage. This step is free of any message. Then replace $\pad$ by an independent uniform value $\pad^\ast$, which changes the game by at most the puncturable-PRF advantage of $\PRF_5$ at $\handle_j$. In the reduction to this advantage, the challenge value plays the role of the pad both inside the circuit (through the hardwired point) and in query $j$'s own salt, so the query-$j$ output, which does depend on the pad, is simulated consistently; the adversary's other evaluations at handle $\handle_j$ take the main branch, since it cannot make the trigger activate there without knowing the pad. The point $\handle_j$ contains no message, so no guess is made. The other signing queries use other handles and are unaffected.

We now show the two procedures produce identically distributed signatures, given the uniform pad. Write the query-$j$ signature as $(P, s, \salt)$. In both procedures the signature verifies, so $s = \mathsf{dlog}_g(P) + \HFEv(\hk, (P, m, \salt)) \cdot x \bmod p$; that is, $s$ is a fixed function of $(P, \salt)$ and the fixed key $x$. It therefore suffices to show that the pair $(P, \salt)$ has the same distribution under the two procedures.

Under the honest procedure, $P = g^r$ is uniform in $\GG$ and $\salt = (\handle, w)$ is uniform in $\bits^{\secp} \times \bits^{\plen}$, and the two are independent.

Under the trigger procedure, $\handle$ is uniform and drawn independently of the payload $(c,z,h,t)$, and $w = (c, z, h, t) \oplus \pad^\ast$ with $\pad^\ast$ uniform and independent of the payload. By the one-time pad, $w$ is uniform and independent of the payload, hence independent of $(c, z)$. So $w$ and $\handle$ are together independent of $(c,z)$, and $P = g^{z} X^{-c}$ is a function of $(c,z)$ alone, so $\salt = (\handle, w)$ is independent of $P$. Finally $P$ is uniform in $\GG$ because $z$ is uniform, and $\salt$ is uniform because $\handle$ and $w$ are. So $(P,\salt)$ is uniform and its two parts are independent, matching the honest procedure.

The two distributions of $(P, \salt)$ agree, so the two distributions of $(P, s, \salt)$ agree. The only slack is the event, of probability at most $2^{-\secp}$, that the honest procedure's uniform salt happens to decode to a payload that activates the trigger; on its complement the honest signature uses the main branch as assumed.
\end{proof}

\paragraph*{Discussion.} The lemma is the point of the encoding. The one-time pad, keyed by a handle that the reduction chooses freshly and that does not mention the message, makes the salt independent of the programmed data $(c, z)$ even though that data determines the signature's point. Verification then fixes the response. So an honest signature and a trigger signature are the same object. Chaining the lemma over the query index $j = 1, \ldots, q$, in a sequence of hybrids where the $j$-th moves query $j$ from honest to trigger, moves every signature to the trigger procedure. Each hop is one application of the lemma, so the total loss is $q$ times the per-hop bound, which is polynomial: there is no message guess, only the message-free handle-hardwiring inside each hop.

Compared with the signing hops of Theorem~\ref{thm-adaptive}, the factor $2^{\msglen}$ is gone, and it is worth saying exactly why, since this is the crux of the improvement. In both the static proof of Theorem~\ref{thm-uf} and the salted proof of Theorem~\ref{thm-adaptive}, the reduction makes the main-branch masks $(a,b)$ uniform at the signature points, by puncturing $\PRF_1, \PRF_2$ there. Puncturing at these points is not itself the difficulty, because they form a set of only $q$ points and puncturable PRFs support puncturing at any polynomial-size set. The difficulty is that each such point contains the message, so to place the puncture the reduction must know the message in advance, and for an adaptive query it guesses the digest. The encoding removes this need entirely. Here the main-branch masks $(a,b)$ are never made uniform and $\PRF_1, \PRF_2$ are never punctured, because the value $\HFEv(\hk, (P,m,\salt))$ never needs to be uniform: the free data $(c,z)$ carried in the salt already makes the signature's point uniform, and verification ties the response to whatever the circuit returns. The only PRF touched in the signing hops is $\PRF_5$, punctured at the message-free handle. In short, the encoding trades a message-dependent puncture of $\PRF_1, \PRF_2$ for a message-free puncture of $\PRF_5$, and that is what turns the guess into a free step.

\subsection{The Remaining Requirement, and an Attack Without It} \label{sec-poly-gap}

After all signing queries are trigger-answered, the game runs without the secret key, and a forgery on a fresh message $m^\ast$ is either a main-branch forgery or a trigger-branch forgery. The main-branch case converts to a CDL solution exactly as in Theorem~\ref{thm-adaptive}, with no guess. The trigger-branch case is the remaining step, and it is where message binding matters. We state what is needed and show that the naive encoding fails without it.

For a trigger-branch forgery, the adversary submits $(R^\ast, s^\ast, (\handle^\ast, w^\ast))$ on $m^\ast$ such that the circuit activates on $(R^\ast, m^\ast, (\handle^\ast, w^\ast))$. Activation requires the decoded payload $(c^\ast, z^\ast, h^\ast, t^\ast) = w^\ast \oplus \PRFEv_5(k_5, \handle^\ast)$ to satisfy $t^\ast = \PRFEv_6(k_6, (\handle^\ast, c^\ast, z^\ast, h^\ast))$ and $h^\ast = \CRHF(m^\ast)$ and $R^\ast = g^{z^\ast} X^{-c^\ast}$. The requirement is that the adversary cannot cause this, except with negligible probability, given the signatures it received.

\paragraph*{The attack when the tag is absent.} Suppose the payload were $(c, z, h)$ with no tag $t$, and the circuit checked only $h = \CRHF(m)$ and $R = g^{z} X^{-c}$. The one-time pad is malleable, and this breaks the scheme. The adversary asks for a signature on $m_1$, receiving $(R_1, z_1, (\handle_1, w_1))$ with $w_1 = (c_1, z_1, \CRHF(m_1)) \oplus \pad_1$ for $\pad_1 = \PRFEv_5(k_5, \handle_1)$. It picks any $m^\ast \ne m_1$ and sets
$$w^\ast \gets w_1 \oplus \big(0,\, 0,\, \CRHF(m_1) \oplus \CRHF(m^\ast)\big) \;.$$
Then $w^\ast \oplus \pad_1 = (c_1, z_1, \CRHF(m^\ast))$, so the circuit on input $(R_1, m^\ast, (\handle_1, w^\ast))$ decodes to a payload whose message field is $\CRHF(m^\ast)$ and whose point is $R_1 = g^{z_1} X^{-c_1}$, and the trigger activates and returns $c_1$. The tuple $(R_1, z_1, (\handle_1, w^\ast))$ is then a valid signature on $m^\ast$, which was never queried. This is a universal forgery. The tag is there to stop exactly this: flipping the message field must invalidate the payload.

\paragraph*{What the tag needs to provide.} With the tag, a trigger forgery requires a valid $t^\ast = \PRFEv_6(k_6, (\handle^\ast, c^\ast, z^\ast, h^\ast))$ for a payload the adversary did not receive verbatim, in particular with $h^\ast = \CRHF(m^\ast)$ for a fresh $m^\ast$. Informally, the adversary must forge an authentication tag for a new message under the key $k_6$. The difficulty is that $k_6$ is inside the obfuscated circuit, so the adversary holds a program that checks tags under $k_6$, and the standard reduction to the unforgeability of $\PRF_6$ punctures $k_6$ at the challenge point $(\handle^\ast, c^\ast, z^\ast, \CRHF(m^\ast))$, which contains $\CRHF(m^\ast)$ and is adaptive. Puncturing there requires knowing $m^\ast$ in advance, which reintroduces the guess this section set out to remove. \authnote{Adam}{This is the crux. The signing hops are exact and message-free (Lemma~\ref{lem-exact-sign}); the whole subexponential loss has been pushed into this one authentication step. Removing the guess here means realizing an adaptively-sound, obfuscation-embeddable authentication of the message field whose soundness reduction does not puncture at a message-dependent point. Candidate directions: (1) a statistically-sound binding, where the tag is replaced by a check that is information-theoretically hard to satisfy off the queried payloads, using the fresh-handle structure; (2) a hidden-triggers/SW-style~\cite{STOC:SahWat14} encoding where the message field is bound through a chain that supports single-point puncturing at a message-free coordinate; (3) importing an adaptively-sound puncturing for this specific comparison-gated use. I believe (1) is the most likely to close, because when the forgery handle $\handle^\ast$ is fresh the pad $\PRFEv_5(k_5, \handle^\ast)$ is unknown to the adversary and the decoded payload is unpredictable, so only the reused-handle case needs the tag; the reused-handle case is where malleability bites, and there the constraint is that the adversary flips a payload it has seen. Whether that can be made statistically sound without an authentication key inside the circuit is the open question.}

\paragraph*{Summary.} The encoding turns the adaptive signing simulation into an exact, message-free argument (Lemma~\ref{lem-exact-sign}), which is the bulk of the proof. What remains is a single message-binding step for trigger-branch forgeries. Under polynomially hard iO and PRFs, closing that step gives adaptive EUF-CMA of the encoded-salt scheme for every polynomial number of queries, without placing a query bound in the scheme; we leave it open, and we regard it as the central technical question left by this work.
